\chapter{问题定义}
\label{chap:problem}

\section{项目背景与 NABCD 分析}
“谁能打开你的窗”关注陌生人文字交流中的开场犹豫、回应含义和互动节奏，以人物发现、短纸条、持续对话和感受记录构成低压力交流闭环。用户自愿展示资料，自行决定回应及发送内容，系统不替用户决定关系。

\textbf{需求（Need）。}面向愿意尝试文字交流且希望保留选择权的课程原型使用者，明确对象是真人账号还是 AI 虚拟人物，并提供介绍、接收偏好和可编辑草稿。\textbf{方法（Approach）。}浏览器与 Node.js 服务端协同，SQLite 按账号保存数据；介绍经同意提取标签，双向评分后加权发现对象，收藏后以纸条开启交流。真人本人回信，虚拟人物由模型回复。\textbf{收益（Benefit）。}接收开关、收藏资格及投递额度限定首次接触；表达建议降低组织语言的操作负担而不自动发言，场景与形象支持个性化，心笺与回望支持个人记录。\textbf{替代方案（Competition）。}相较已有联系人即时通信、纯虚拟聊天或推荐列表，本项目展示发现、准入和对话的完整链路，仅作定性方案比较。\textbf{交付（Delivery）。}交付本地浏览器 Demo、源码与课程报告，按“注册与资料—遇见与收藏—纸条与回信—对话与建议—场景与形象”演示。

上述场景与收益用于指导原型设计，后续可通过访谈、试用记录和统一量表评估用户体验。

\section{项目基本目标}
目标是在本地可复现一个规则清晰、状态可恢复的交流原型：建立注册、登录、退出与单账号单有效会话；实现资料标签、双向发现和收藏；以纸条准入建立真人对话并区分虚拟模型互动；提供可编辑建议、情绪参考及失败回退；分别管理头像、背景和立绘，生图经授权、预览和确认才保存；演示虚拟心笺、当期意愿及周期揭晓。验收关注权限、状态转移、持久化、过期结果绑定和重试，完整规约以项目需求文档为准。

\textbf{边界与伦理。}匹配百分数是标签契合度而非成功概率，意图、情绪和交流参考分不确认真实心理，模型自评置信度不是校准概率。Demo 来信、日期推进、意愿和揭晓仅属虚拟状态，AI 不冒充真人回信；未成年虚拟角色应保持适龄日常与友谊交流，人设提示不等于完备保护机制。心笺不发给文本模型，但服务端数据库未加密，不承诺管理员不可见或端到端隐私。密码找回、邮箱及真实身份或年龄核验、真实举报审核和公网生产保障未完成。

\section{可行性与交付条件}
\textbf{技术与操作。}原生 HTML、CSS、JavaScript ES Modules 和 Fetch 配合 Node.js 24 内置 HTTP、Crypto、SQLite 模块，无第三方 npm 业务依赖；服务端校验有限动作、账号归属及事务约束，不接受客户端整份状态覆盖。Windows 启停和状态工具支持本机演示，导航覆盖登录、遇见、纸条、对话与“我的”，提供表单提示、动画跳过及二次确认。手机、桌面与减少动态效果是适配目标，当前版本仍需浏览器验证。

\textbf{模型与资源。}文本采用兼容 Chat Completions 的协议，图片编辑采用携参考图的 Multipart 请求；凭据仅在服务端，返回结构及范围需校验。可用性受配置、网络、协议和额度影响，真人已有会话可在模型故障时手动发送，虚拟待回复则须重试。并发、次数及超时控制资源，不代表实测成本或承载人数。预算\underline{\hspace{1.4cm}}、模型实际费用\underline{\hspace{1.4cm}}、总工时\underline{\hspace{1.4cm}}均待真实资料填写。

\textbf{风险与证据。}默认仅监听本机，可信局域网需配置；HTTP 与数据库均未加密，同步数据库及密码派生可能阻塞高负载。模型迟到、任务重启中断和图片偏差通过超时、消息绑定、版本检查及预览确认降低风险；普通上传检查声明、签名与大小，生成 PNG 才另查解码和透明通道。源码确认能力，\texttt{VERIFICATION.md} 记录带日期的验证，截图仅说明采集时呈现。2026年10月31日120项回归使用模拟模型，非本次重跑、真实模型质量或全面安全审计。

\section{人员与进度管理}
成员信息、实际分工、日期和工时待填写。下表提供建议责任范围，按实际人数调整，并关联任务记录、提交与验证结果。
\begin{table}[H]
  \centering
  \caption{团队信息与责任记录（个人资料待填）}
  \label{tab:problem-members}
  {\songti\rmfamily\zihao{5}
  \renewcommand{\arraystretch}{1.2}
  \begin{tabular}{@{}p{.08\textwidth}p{.12\textwidth}p{.14\textwidth}p{.32\textwidth}p{.20\textwidth}@{}}
    \toprule
    序号 & 姓名 & 学号 & 拟记录的责任范围 & 实际贡献证据 \\
    \midrule
    1 & 待填 & 待填 & 需求梳理、业务规则与报告 & 待填 \\
    2 & 待填 & 待填 & 前端交互、场景与形象界面 & 待填 \\
    3 & 待填 & 待填 & 服务端认证、数据与通信 & 待填 \\
    4 & 待填 & 待填 & 模型接入、验证与交付 & 待填 \\
    \bottomrule
  \end{tabular}}
\end{table}
阶段顺序按任务依赖归纳，计划与实际日期待补。每阶段记录负责人、输入输出、问题处理及验证结果。
\begin{table}[H]
  \centering
  \caption{阶段交付与进度核对表}
  \label{tab:problem-milestones}
  {\songti\rmfamily\zihao{5}
  \renewcommand{\arraystretch}{1.2}
  \begin{tabular}{@{}p{.19\textwidth}p{.45\textwidth}p{.12\textwidth}p{.14\textwidth}@{}}
    \toprule
    阶段 & 交付与核对内容 & 计划时间 & 实际时间 \\
    \midrule
    需求与边界 & 明确真人、虚拟人物、授权与未实现能力 & 待填 & 待填 \\
    基础设施 & 账号、会话、数据库及动作校验 & 待填 & 待填 \\
    发现与通信 & 标签、抽取、收藏、纸条和真人对话 & 待填 & 待填 \\
    AI 与媒体 & 表达建议、情绪分析、媒体与生图任务 & 待填 & 待填 \\
    验证与交付 & 隔离测试、截图核对、限制说明及报告 & 待填 & 待填 \\
    \bottomrule
  \end{tabular}}
\end{table}
